% ---------------------------------------------------------------------
% Discussion, Limitations, and Ethics (~1 page).
% ---------------------------------------------------------------------

\subsection{What the Results Mean for Practice}

Three fundamental practical insights emerge from the empirical evaluations in Section~\ref{sec:results}:

First, \emph{in-domain benchmark leaderboards fail to predict real-world operational reliability}: models that achieve state-of-the-art in-domain EER on ASVspoof~2019~LA (down to 2.86\% for B5 and 6.81\% for AuralGuard v2) diverge dramatically under zero-shot domain shift, spanning 16.55\% to 50.62\% EER on In-the-Wild (Table~\ref{tab:main}). Traditional spectral-temporal architectures such as LFCC-LCNN (50.62\% EER) and RawNet2 (38.08\% EER) suffer catastrophic performance degradation when confronted with unconstrained recording channels and varied speaker acoustics. Procurement or forensic deployment decisions based exclusively on clean laboratory benchmarks are fundamentally misleading.

Second, \emph{foundation-model features provide indispensable contextual invariance but do not eliminate vocoder vulnerabilities}: pre-trained representations from WavLM-Large prevent the total collapse observed in raw acoustic models on found audio (B5 and AuralGuard maintain $<$20\% EER on In-the-Wild compared to $>$38\% for non-SSL systems). However, on the multi-vocoder WaveFake benchmark, all evaluated systems—including single-view foundation models—exhibit substantial error rates (45--50\% EER). This confirms that standard SSL pre-training objectives (masked speech prediction and denoising) compress high-frequency phase and fine vocoder synthesis artifacts, which must be reinforced via specialized physical artifact pathways~\cite{singh2026resilient,nagar2026hybrid}.

Third, \emph{decision calibration must be treated as a first-class evaluation dimension}: raw classification confidence is treacherous in forensic contexts. While legacy architectures like AASIST appear well-calibrated in-domain ($\mathrm{ECE} = 0.0545$), their calibration degrades catastrophically out-of-domain ($\mathrm{ECE} = 0.5552$ on In-the-Wild, Table~\ref{tab:calibration}). AuralGuard's multi-view one-class architecture maintains superior calibration under severe distribution shift ($\mathrm{ECE} = 0.0965$ for v1, $0.1716$ for v2). Operational forensic systems must continuously monitor posterior calibration on unconstrained traffic rather than relying on in-domain probability assertions.

\subsection{Forensic Admissibility, Regulatory Compliance, and the Calibration Imperative}
\label{ssec:forensic_legal}

The deployment of audio deepfake detection systems in judicial proceedings, criminal investigations, and regulatory oversight introduces stringent evidentiary standards that conventional classification accuracy metrics do not satisfy. Under the landmark \emph{Daubert} standard (\emph{Daubert v. Merrell Dow Pharmaceuticals, Inc.}~\cite{daubert1993}) and Federal Rules of Evidence Rule 702, scientific expert testimony based on computational techniques requires: (i)~an established, empirically validated error rate under operational conditions, (ii)~peer-reviewed methodology, and (iii)~known standards of operational control. 

Concurrently, Article 50 of the European Union Artificial Intelligence Act (Regulation EU 2024/1689)~\cite{eu_ai_act_2024} imposes explicit transparency and disclosure obligations on deployers of AI systems that generate or manipulate audio content, mandating that deepfakes be marked and detectable with verifiable fidelity. In forensic practice, presenting a court with an uncalibrated countermeasure that outputs a ``99.4\% probability of synthesis'' on a recording—when the system's empirical accuracy under that confidence bin is merely 52\% (as demonstrated by legacy baselines in Figure~\ref{fig:reliability})—severely compromises evidentiary integrity and risks wrongful conviction or false exoneration.

AuralGuard directly satisfies these evidentiary and regulatory imperatives through two design principles:
\begin{enumerate}
\item \textbf{Calibrated Probability Estimation:} By anchoring representations to a compact bona fide centroid $\mathbf{w}_0$ and constraining score distributions via angular margin penalization, AuralGuard achieves an out-of-domain ECE of $0.0965$ (v1) on unconstrained audio. As shown in the reliability diagram (Figure~\ref{fig:reliability}), predicted posterior probabilities track observed empirical frequencies across the entire probability continuum, enabling forensic examiners to offer probability statements that genuinely reflect statistical likelihood.
\item \textbf{Statistical Uncertainty Quantification:} Rather than presenting isolated point estimates, our framework provides non-parametric bootstrap confidence intervals (Table~\ref{tab:main} and Figure~\ref{fig:forest}) across $B=1,000$ resamples. This allows expert witnesses to report rigorous confidence intervals (e.g., In-the-Wild EER: $[18.78\%, 19.61\%]$ for v2) required to withstand judicial scrutiny.
\end{enumerate}

\subsection{Threat Model and Robustness to Multi-Stage Adversarial Laundering}
\label{ssec:laundering}

In real-world disinformation campaigns, synthetic audio rarely arrives in pristine uncompressed WAV format. Instead, malicious actors deploy \emph{adversarial laundering pipelines} to scrub telltale synthesis artifacts before distributing deepfakes across social platforms:
\begin{itemize}
\item \textbf{Lossy Perceptual Transcoding:} Audio is uploaded to encrypted messaging channels (e.g., WhatsApp or Telegram), subjected to Opus transcoding at low bitrates (16--32~kbps), downloaded, and subsequently shared on platforms such as TikTok or YouTube Shorts using lossy AAC or MP3 encoding (64--128~kbps). The successive psychoacoustic quantization steps eliminate fine high-frequency phase signatures and discard sub-band energy above 8~kHz.
\item \textbf{Acoustic Re-Recording (Air-Chamber Splicing):} Perpetrators replay synthetic audio through low-fidelity loudspeakers in reverberant rooms and re-record the audio via smartphone microphones, superimposing physical room impulse responses (RIR) and ambient acoustic noise to mask vocoder periodicity markers.
\end{itemize}

A single-stream acoustic detector relying exclusively on high-frequency Fourier phase or cepstral coefficients fails catastrophically under such laundering pipelines, as demonstrated by the collapse of LFCC-LCNN on In-the-Wild. Conversely, single-stream foundation models relying exclusively on semantic SSL representations can be fooled by high-fidelity neural vocoders that preserve natural prosody. AuralGuard's bidirectional gated cross-attention (Section~\ref{ssec:fusion}) provides intrinsic resilience against multi-stage laundering: when lossy transcoding obliterates phase dynamics, the gating mechanism downweights corrupted artifact channels and queries the robust foundation-model representations. Conversely, when semantic prosody is mimicked by large generative text-to-speech models, the artifact branch captures surviving sub-band phase jitter and group-delay anomalies.

\subsection{Specialized Detectors versus General Audio-Language Models}
\label{ssec:specialized_vs_allm}

Our findings contribute to an ongoing debate in multimedia forensics: whether general-purpose audio-language foundation models (ALLMs) will render specialized countermeasures obsolete~\cite{gong2025allm4add}. Recent investigations demonstrate that instruction-tuned audio foundation models prompted zero-shot operate near chance on speech anti-spoofing tasks~\cite{gong2025allm4add,oyucu2026generalization}. This outcome is architectural rather than accidental: foundation models and neural audio codecs are explicitly optimized to compress semantic linguistic content while discarding phase irregularities, spectral ripples, and high-frequency quantization noise—the exact physical cues that betray synthetic speech. Consequently, specialized countermeasures such as AuralGuard remain essential, providing calibrated artifact-level anomaly detection that can serve as trusted evidence streams for upstream multimodal reasoning systems.

\subsection{Limitations}
\label{ssec:limitations}

Despite the empirical gains documented in Section~\ref{sec:results}, several technical limitations warrant candid discussion:
\begin{enumerate}
\item \textbf{White-Box Adaptive Adversaries:} Our evaluation validates robustness against unseen generative architectures and unconstrained acoustic environments, but does not evaluate an adversary with full gradient access to AuralGuard who optimizes anti-forensic perturbation vectors constrained by psychoacoustic masking thresholds. Developing certified robustness bounds and adversarial training curricula remains an active frontier.
\item \textbf{Computational Footprint for Edge Devices:} Although the WavLM-Large backbone remains frozen throughout training, its 315M parameters and 71.4~GMACs inference footprint impose hardware constraints. While our measured CPU real-time factor ($\mathrm{RTF} = 0.412$) comfortably sustains high-throughput server-side verification, real-time deployment on embedded edge devices (e.g., smart voice assistants or mobile transceivers) will require structured knowledge distillation into compressed transformer back-ends~\cite{liu2025nes2net,nagar2026hybrid}.
\item \textbf{Vocoder Generalization Bottleneck:} The performance plateau observed on WaveFake (45.19\% EER for v2) reveals that zero-shot transfer across structurally distinct neural vocoders (e.g., HiFi-GAN vs. MelGAN vs. Parallel WaveGAN) remains challenging when training data is restricted to ASVspoof~2019~LA. Future iterations must incorporate vocoder-diverse synthetic pre-training.
\item \textbf{Partial Temporal Spoofing:} The current AuralGuard formulation evaluates whole-utterance authenticity. In targeted disinformation scenarios involving selective word substitution or temporal splicing~\cite{luong2024llamapartialspoof}, frame-level attention pooling must be extended to localize and segment spliced temporal intervals.
\end{enumerate}

\subsection{Ethical Considerations and Dual-Use Governance}
\label{ssec:ethics}

Audio deepfake detection is inherently dual-use. AuralGuard is designed to mitigate harm from voice-cloning financial fraud, non-consensual biometric impersonation, and political audio disinformation. However, forensic classification systems can also introduce acute risks if misapplied: an overconfident, miscalibrated false-alarm verdict can wrongly discredit authentic whistleblower recordings or genuine legal evidence—a reverse harm documented in recent investigative journalism~\cite{chandra2025deepfakeeval}. 

To counteract these harms, AuralGuard implements three explicit safeguards:
(i)~it outputs calibrated posterior probabilities with non-parametric uncertainty bounds rather than arbitrary binary verdicts,
(ii)~it exposes internal evidence attribution via gate activation statistics, and
(iii)~it is explicitly designated as an evidentiary triage tool that must corroborate, rather than replace, human forensic deliberation. 

All training and evaluation corpora utilized in this research are publicly available academic datasets collected under established institutional review board (IRB) protocols and open research licenses; no private voice data was scraped or utilized without consent. Complete evaluation protocols and model checkpoints are made available to foster reproducible defensive research, while weaponizable adaptive evasion tools are withheld.
